% ==========================================================================
% GENERATED by scripts/split_paper.py. Do not edit this file.
%
% Source: imported_papers/grokking/iclr2027_conference_v2.tex
% Part 13: Theory of credit equalization
%
% Edits belong upstream, in the published paper. To change what the
% chapter includes, change imports/grokking.toml and re-run
% the splitter.
% ==========================================================================
\section{Theory of credit equalization}
\label{app:theory}
We use a diagonal linear network to examine how equalizing gradient magnitudes changes the distribution of contributions across units. The model allows an exact comparison between unmodified gradient flow, per-unit equalization and groupwise equalization, and it explains a mechanism for the redistribution; the ReLU-network experiments test the effect itself. We first record two elementary properties of the participation ratio and of the intervention.

\begin{proposition}[Participation ratio bounds]
\label{prop:pr}
For $a \in \R^n_{\geq 0}$, $a \neq 0$, the participation ratio of Equation~\ref{eq:pr} satisfies $1/n \leq \mathrm{PR}(a) \leq 1$. The upper bound is attained exactly when all $a_j$ are equal, and the lower bound exactly when a single $a_j$ is nonzero. Consequently $n \cdot \mathrm{PR}(a)$ interpolates between $1$ and $n$ and acts as an effective count of contributing units.
\end{proposition}
\begin{proof}
By Cauchy--Schwarz, $(\sum_j a_j)^2 \leq n \sum_j a_j^2$ with equality exactly for constant $a$, which gives $\mathrm{PR} \leq 1$. Since the $a_j$ are nonnegative, $(\sum_j a_j)^2 \geq \sum_j a_j^2$ with equality exactly when at most one entry is nonzero, which gives $\mathrm{PR} \geq 1/n$.
\end{proof}

\begin{proposition}[Credit rescaling preserves a descent direction]
\label{prop:descent}
Let $g \neq 0$ be the loss gradient of the weight matrix to which Equation~\ref{eq:knob} is applied, with one row $g_j$ per neuron of the layer it feeds, and let $\tilde{g}$ be the rescaled signal for any $\beta \in [0,1]$. Then $\langle \tilde{g}, g \rangle > 0$, so a sufficiently small step along $-\tilde{g}$ decreases the loss when the other parameters are held fixed.
\end{proposition}
\begin{proof}
Each nonzero row is rescaled by the positive scalar $s_j = \big((1-\beta)\|g_j\| + \beta \bar{g}\big) / \|g_j\|$, and zero rows remain zero, so
$\langle \tilde{g}, g \rangle = \sum_{j:\, \|g_j\| > 0} s_j \|g_j\|^2 > 0$ whenever $g \neq 0$. The directional derivative of the loss along $-\tilde{g}$ is therefore negative.
\end{proof}
This statement concerns the rescaled loss gradient before Adam and weight decay.

\begin{proposition}[Unmodified and equalized flow select different minima]
\label{prop:select}
Consider a diagonal linear network with $n$ units. Unit $i$ has incoming weight $u_i$, outgoing weight $v_i$ and contribution $a_i = u_i v_i$. Initialize $u_i = v_i = \sqrt{\epsilon_i} > 0$ with $\sum_i \epsilon_i < 1$, and train by gradient flow on the loss
\begin{equation}
L \;=\; \tfrac{1}{2}\Big(\sum\nolimits_i a_i - 1\Big)^{2},
\label{eq:toyloss}
\end{equation}
whose minima form the set $\{a : \sum_i a_i = 1\}$. In this model, equalization rescales the two-coordinate gradient $(\partial_{u_i} L, \partial_{v_i} L)$ of each unit by one common factor, so both weights of a unit are updated symmetrically and stay balanced, $u_i = v_i$ and $a_i = u_i^2$. Participation is computed from the contributions $a_i$. Unless the $\epsilon_i$ are all equal, unmodified and equalized gradient flow converge to different minima:
\begin{itemize}
\item[(i)] Unmodified gradient flow keeps every ratio $a_i / a_j$ at its initial value, so it converges to the rescaled initialization $a_i = \epsilon_i / \sum_j \epsilon_j$, whose participation ratio equals that of the initialization.
\item[(ii)] The symmetric per-unit equalized flow keeps every difference $u_i - u_j$ at its initial value, so it converges to $u_i = \sqrt{\epsilon_i} + C$ with $C$ set by $\sum_i (\sqrt{\epsilon_i} + C)^2 = 1$.
\item[(iii)] $\mathrm{PR}\big((\sqrt{\epsilon_i} + C)^2\big)$ is nondecreasing in $C \geq 0$, strictly increasing unless the $\epsilon_i$ are equal. Whenever the initialization undershoots the constraint, $\sum_i \epsilon_i < 1$ and $C > 0$, the equalized solution has strictly higher participation than the initialization, and it reaches $1$ as the initialization scale tends to zero. Overshooting initializations are outside this statement because some coordinates can cross zero under equalization.
\end{itemize}
\end{proposition}
\begin{proof}
Write $r = \sum_j a_j - 1$ for the residual, a single scalar shared by every unit. The two flows differ in one line,
\begin{equation}
\text{unmodified flow: }\ \dot{u}_i = -\partial_{u_i} L = -r\,u_i,
\qquad\qquad
\text{equalized flow: }\ \dot{u}_i = -r\,\bar{u},
\quad \bar{u} = \tfrac{1}{n}\textstyle\sum_j u_j,
\label{eq:twoflows}
\end{equation}
where the equalized form holds because unit $i$'s two-coordinate gradient $(r u_i,\, r u_i)$ has norm $|r|\sqrt{2}\,u_i$, and equalization rescales it to the mean norm $|r|\sqrt{2}\,\bar{u}$. For (i), $\dot{a}_i = 2 u_i \dot{u}_i = -2 r a_i$, so $\frac{d}{dt} \log(a_i / a_j) = 0$ because $r$ is common to all units. The invariance is a property of this model, since in a network whose units see different residuals the ratios do move. For (ii), every $u_i$ takes the same step, so $\sqrt{a_i(t)} = \sqrt{\epsilon_i} + C(t)$ with $C(t) = -\int_0^t r \bar{u}\, ds$, and fitting the target fixes $C$. For (iii), write $b_i = \sqrt{\epsilon_i} + C$ and $S_k = \sum_i b_i^k$, so that $\mathrm{PR} = S_2^2 / (n S_4)$ and
\begin{equation}
\frac{d}{dC}\,\mathrm{PR} \;=\; \frac{4\,S_2\,\big(S_1 S_4 - S_2 S_3\big)}{n\,S_4^{2}},
\qquad
S_1 S_4 - S_2 S_3 \;=\; \tfrac{1}{2}\sum_{i,j} b_i b_j\,(b_i + b_j)(b_i - b_j)^2 \;\geq\; 0,
\label{eq:prmono}
\end{equation}
where the identity follows by symmetrizing the double sum $\sum_{i,j} b_i b_j^4 - b_i^2 b_j^3$ over $i$ and $j$, and it is zero exactly when all nonzero $b_i$ coincide. Participation is nondecreasing along the path from $C = 0$, where it equals $\mathrm{PR}(\epsilon)$, to the fitted $C > 0$, and strictly increasing unless the $\epsilon_i$ are equal. As the initialization scale tends to zero, $C \to 1/\sqrt{n}$ and $a_i \to 1/n$, whose participation ratio is $1$.
\end{proof}

Under unmodified gradient flow every contribution grows by the same factor, so their ratios stay fixed. Under symmetric per-unit equalization every $u_i$ receives the same increment; the initial differences then become small relative to the shared growth, and participation increases. Both flows optimize the same loss but select different members of its set of minima. The model differs from the ReLU experiment in two ways: it rescales both weights of a unit so that they stay balanced, whereas the experiment rescales only the incoming weights, and it measures the products $u_i v_i$ rather than the outgoing norms. It provides a mechanism for the observed redistribution, not an exact prediction for the trained nonlinear network. The same model describes groupwise equalization, which preserves the ratios within each group while equalizing the group totals.

\begin{proposition}[Group equalization removes only between-group inequality]
\label{prop:grain}
Partition the $n$ units into $g$ groups of $m = n/g$, fixed throughout training. In the model of Proposition~\ref{prop:select}, rescale each unit's two-coordinate gradient by the factor, common within its group, that equalizes the total credit across groups. Write
\begin{equation}
p_i \;=\; \frac{\sqrt{\epsilon_i}}{\sum_{j \in \mathrm{group}(i)} \sqrt{\epsilon_j}}
\label{eq:groupshare}
\end{equation}
for unit $i$'s share of its own group's initialization, measured in the $\sqrt{a}$ coordinate in which the flow is linear, and let $k$ run over groups. As the initialization scale tends to zero, the fitted solution has participation ratio
\begin{equation}
\mathrm{PR}(m) \;=\; \frac{S^2}{m\,Q},
\qquad
S = \frac{1}{g}\sum_{k}\sum_{i \in k} p_i^2,
\qquad
Q = \frac{1}{g}\sum_{k}\sum_{i \in k} p_i^4,
\label{eq:grainlaw}
\end{equation}
which equals the unmodified-flow value of Proposition~\ref{prop:select} at $g=1$ and equals $1$ at $m=1$.
\end{proposition}

\begin{proof}
The rescaling factor is shared within a group, so inside group $k$ every $\log u_i$ has the same time derivative. The within-group shares $p_i$ freeze at their initial values, as the ratios freeze under unmodified flow. The group totals $U_k = \sum_{i \in k} u_i$ obey
\begin{equation}
\dot{U}_k \;=\; -\,r\,c_k U_k \;=\; -\,r\,\bar{U},
\qquad c_k = \bar{U}/U_k,
\qquad \bar{U} = \tfrac{1}{g}\textstyle\sum_k U_k,
\end{equation}
a common increment, so as the initialization tends to zero all group totals equalize to a common $U$. The solution is then $a_i = p_i^2 U^2$ with $U^2 = 1/(gS)$ forced by $\sum_i a_i = 1$, and substituting into Equation~\ref{eq:pr} gives Equation~\ref{eq:grainlaw}. At $m=1$ every share is $1$, so $S = Q = 1$ and $\mathrm{PR} = 1$, recovering Proposition~\ref{prop:select}. At $g=1$ the rescaling is the identity, and Equation~\ref{eq:grainlaw} reduces to the participation of the normalized initialization, the unmodified-flow value.
\end{proof}

Groupwise equalization changes the totals assigned to the groups while preserving the relative weights inside each group. To compare Equation~\ref{eq:grainlaw} with the neural-network sweep, we assume lognormally distributed initial contributions with log standard deviation $\sigma$. For many units the unmodified model then has participation $e^{-\sigma^2}$; we set this equal to the measured whole-layer value $0.602$, which fixes $\sigma$ from that one point. The model then predicts participation $0.674$, $0.737$, $0.834$ and $1.000$ for groups of eight, four, two and one, against measured values of $0.623$, $0.695$, $0.753$ and $0.992$. It reproduces the shape of the curve: flat for large groups and steep at the smallest.

\begin{proposition}[Participation and score fluctuations under random deletion]
\label{prop:fluct}
Model the score of the correct class as a sum of unit contributions, $y = \sum_j c_j$ with $c_j \geq 0$ over $n$ units, and delete units independently with probability $\rho$, giving the surviving score $y_\rho = \sum_j b_j c_j$ with $b_j \sim \mathrm{Bernoulli}(1-\rho)$ marking survival. Then, assuming $\sum_j c_j>0$ and $0\leq \rho<1$, for every $\varepsilon \in (0,1)$,
\begin{equation}
\frac{\sqrt{\mathrm{Var}[y_\rho]}}{\mathbb{E}[y_\rho]}
= \sqrt{\frac{\rho}{1-\rho}} \cdot \frac{1}{\sqrt{n \cdot \mathrm{PR}(c)}},
\qquad
\mathbb{P}\big(y_\rho < (1-\varepsilon)\,\mathbb{E}[y_\rho]\big)
\;\leq\; \frac{\rho}{(1-\rho)\,\varepsilon^2 \, n \cdot \mathrm{PR}(c)},
\label{eq:fluct}
\end{equation}
where $\mathrm{PR}(c)$ is Equation~\ref{eq:pr} evaluated on the vector of contributions.
\end{proposition}
\begin{proof}
$\mathbb{E}[y_\rho] = (1-\rho) \sum_j c_j$ and, by independence, $\mathrm{Var}[y_\rho] = \rho(1-\rho) \sum_j c_j^2$. Dividing and using Equation~\ref{eq:pr} gives the first identity, and Chebyshev's inequality gives the tail bound.
\end{proof}

In this additive model, a larger effective count reduces the relative fluctuation of the score under independent random deletion. The proposition gives the direction of the effect; the deletion experiments establish the accuracy differences in the trained networks.
